\documentclass[tikz,border=8pt]{standalone}
\usepackage{tikz}
\usepackage{amsmath}
\usepackage{pgfplots}
\pgfplotsset{compat=1.18}
\usepackage{ctex}
\usetikzlibrary{calc, positioning, arrows.meta, shapes, fit, backgrounds}

\begin{document}
\begin{tikzpicture}[font=\normalsize]

% 两个相交圆的圆心
\def\cx{1.8}   % 圆心X偏移（左圆在 0, 右圆在 2*cx）
\def\r{2.2}    % 半径

% 左圆（蓝色，代表 H(X)）
\draw[blue!70!black, very thick, fill=blue!10]
  (0,0) circle (\r cm);

% 右圆（红色，代表 H(Y)）
\draw[red!70!black, very thick, fill=red!10]
  (3.6,0) circle (\r cm);

% 交叠区域（互信息 I(X;Y)）——重叠区域着色
\begin{scope}
  \clip (0,0) circle (\r cm);
  \fill[violet!30] (3.6,0) circle (\r cm);
\end{scope}

% 边框重画（让交叠后圆圈清晰）
\draw[blue!70!black, very thick] (0,0) circle (\r cm);
\draw[red!70!black, very thick] (3.6,0) circle (\r cm);

% 区域标注
% H(X|Y)：左独有区域
\node[font=\small, blue!80!black, align=center] at (-1.3, 0)
  {$H(X|Y)$};

% I(X;Y)：交叠区域
\node[font=\small, violet!80!black, align=center] at (1.8, 0)
  {$I(X;Y)$};

% H(Y|X)：右独有区域
\node[font=\small, red!80!black, align=center] at (4.9, 0)
  {$H(Y|X)$};

% 圆的标签（上方）
\node[font=\normalsize\bfseries, blue!80!black] at (-0.6, 2.5) {$H(X)$};
\node[font=\normalsize\bfseries, red!80!black] at (4.2, 2.5) {$H(Y)$};

% H(X,Y) 大括号标注（横跨两圆）
\draw[<->, thick, gray!70] (-2.25, -3.0) -- (5.85, -3.0);
\node[font=\small, gray!70!black] at (1.8, -3.35) {$H(X,Y)$（联合熵）};

% 竖向连线
\draw[gray!50, thin, dashed] (-2.25, 0) -- (-2.25, -3.0);
\draw[gray!50, thin, dashed] (5.85, 0) -- (5.85, -3.0);

% 整体标题
\node[font=\large\bfseries] at (1.8, 3.5) {信息论维恩图};

% 右侧公式框
\node[font=\small, align=left,
      fill=white, draw=gray!50, thin, rounded corners=2pt, inner sep=6pt]
  at (9.0, 0.8) {
    $I(X;Y)=H(X)+H(Y)-H(X,Y)$\\[4pt]
    $H(X,Y)=H(X|Y)+I(X;Y)+H(Y|X)$\\[4pt]
    $I(X;Y)=H(X)-H(X|Y)$\\[4pt]
    $I(X;Y) \geq 0$
  };

\end{tikzpicture}
\end{document}
